% Complete formal comparison and its three-object restriction.
\section{次数別完備な曲率付き移送と形式比較の構成}
\label{cft:sec:transport}

本節では、完備係数上のMF比較を実際の関手と逆関手によって構成し、
その固定三対象への制限を明らかにする。
既知のKontsevich formalityを\emph{曲率を入れる前}に使い、
その後のゲージ経路は有限Artin次数ごとに直接積分する。
曲率の大きさ、全Hochschild cochain、単位、逆写像、有限rank、
二つの次数、指定された周期比較をそれぞれ確認する。
完備化と一般点への局所化の規約は第\ref{gdr:subsec:scope}節に従う。

\subsection{主張と既知結果の使用範囲}
$(R,\mathfrak m)$を剰余体$\mathbb C$の完備noetherian局所$\mathbb C$-代数とする。
ここでは$R$の正則性は不要である。
$S=\mathbb C[x_1,\ldots,x_N]$、$N\ge3$、$\deg x_i=1$とし、
$A^{\rm gr}=\bigoplus_{j\ge0}A_j$をunitalな中央$R$-代数とする。
各$A_j$は有限自由、$A_0=R$、$A^{\rm gr}/\mathfrak m A^{\rm gr}=S$と仮定する。
$\Omega\in Z(A^{\rm gr})_3$の剰余$F$は滑らかな射影三次超曲面を定めるものとする。
次数別に剰余の単項式基底を持ち上げ、単位を保つ加群同定
$A_j\simeq S_j\otimes_{\mathbb C}R$を選ぶ。

\begin{theorem}[完備係数上の実際のgraded MF比較]
\label{cft:thm:complete}
上の入力に対し、$W\in R[x_1,\ldots,x_N]_3$、$W\bmod\mathfrak m=F$と、
有限rank graded MFの実際の$R$-線形unital $A_\infty$同値
\begin{equation}\label{cft:eq:equivalence}
 \operatorname{MF}^{\mathbb Z}(A^{\rm gr},\Omega)
 \simeq\operatorname{MF}^{\mathbb Z}(R[x_1,\ldots,x_N],W)
\end{equation}
が存在する。これは内部shiftと、中心元による$(3)\simeq[2]$を保ち、
剰余では恒等であり、すべての$R/\mathfrak m^{q+1}$への切断と両立する。
有限自由模型上で互いに逆な関手を構成し、その二つのrankと内部shiftを保つ。
構成は$\deg t=(2,-3)$の周期変数についてrelativeであり、$t$を微分しない。
同値は有限coneと冪等完備化へ延長する。

曲率付きbar写像のTaylor成分$F_r$と逆写像の成分$H_r$は
\begin{equation}\label{cft:eq:bound}
 F_0=H_0=0,\quad F_1,H_1=1+O(\mathfrak m),\quad
 F_{2r}=H_{2r}=0,\quad F_{2r+1},H_{2r+1}\in\mathfrak m^r\quad(r\ge1)
\end{equation}
を満たすように取れる。次数付き代数の厳密同型は、この写像の前後に合成する。
この評価はMF行列を挿入する前のbar写像について述べており、
挿入後のMF関手の偶数Taylor成分の消滅を意味しない。
\end{theorem}

Kontsevichの既知結果は、affine formality、HKR、Maurer--Cartan変形問題の比較
\cite[Theorem~4.6, Section~4.6.2, Theorem~6.3]{PKontsevich}である。
曲率付きtwistによるgraded MFの記述には
\cite[Section~2, 特にRemark~2.8とTheorem~2.12]{CT}を使う。
後者を形式的な無限arity関手へ無条件に適用せず、以下で評価式の有限性と収束を証明する。
これらの一般理論の新規性は主張せず、本定理の組合せの初出は未確認である。

\begin{remark}[既知の曲率付きformalityとの比較]\label{cft:rem:prior-formality}
Carqueville--Kay\cite[Proposition~2.3, Remark~2.4, Proposition~3.4]{CK}の
Proposition~3.4は、$S=\mathbb C[x_1,\ldots,x_N]$上の曲率付き代数と
$[-F,-]$を微分とするpolyvector場の変形複体を比較する。
用いられる第二種cochainはarityについての直和であり、各arityでは
全多重線形写像を許すので、変形積がbidifferentialであることを入力条件とはしない。
同論文のTheorem~3.1はさらに固定したMFのendomorphism代数を扱う。
本稿の全graded MF圏の同値や指定周期比較は、同定理の主張そのものではない。

本稿の状況には、$\deg t=(2,-3)$の周期変数を係数として扱い、$t$を微分しない
relativeな比較、二重次数とnormalizedな単位を保つ構成、
$\mathfrak m$進完備化を加える必要がある。
第二種cochainをこの位相で完備化すると、各$R/\mathfrak m^{q+1}$では有限arityとなる。
これは初めからarityの任意の無限直積を許すこととは異なる。
$t$-relativeなpolyvector複体では、勾配Koszul複体により
\begin{equation}\label{cft:eq:prior-jacobian}
 H\bigl(T_{\rm poly}\otimes_{\mathbb C}\mathbb C[t,t^{-1}],[-tF,-]\bigr)
 \simeq\operatorname{Jac}(F)\otimes_{\mathbb C}\mathbb C[t,t^{-1}].
\end{equation}
右辺の関数のLie次数は$-1$であるため、Lie次数一・内部weight零の成分は
$t\operatorname{Jac}(F)_3$である。関数の斉次代表同士のSchouten括弧は零であり、
基点$b_F=\mu_0-tF$の曲率を微分に含めて
$\delta=(\mu_A-\mu_0)-t(\Omega-F)\in\mathfrak m$を扱うことが、
potentialによる形式的存在原理に対応する。

一般の孤立potentialでは、形式係数の極限は多項式次数について有界とは限らず、
例えば$\mathbb C[x][[h]]$と$\mathbb C[[h]][x]$は異なる。
従って、形式的なMC同値だけで本稿の次数別直和上の有限rank同値を結論しない。
以下では補題\ref{cft:lem:model}で元の全cochainから始め、
補題\ref{cft:lem:arity}、\ref{cft:lem:flow}で単位・二重次数と有限評価を確認し、
式\eqref{cft:eq:two-block-matrix}で固定rank・内部shiftのMFに留まることを証明する。
三次特有の定数trivector、一段のゲージ経路、式\eqref{cft:eq:bound}の一様評価は、
この一般原理の本稿での明示的な実現である。
通常のMF圏のHochschild複体と第二種複体を無条件に同一視する必要はない。
\end{remark}

\subsection{通常の結合変形と中心元の有限次元逆像}
\begin{lemma}[次数を保つunitalな星積模型]\label{cft:lem:model}
formal Poisson tensor
$\pi\in\mathfrak m(S_2\otimes\bigwedge^2\langle\partial_i\rangle)\widehat\otimes R$
と、単位・内部次数を保つ厳密な代数同型
\begin{equation}\label{cft:eq:phi}
 \phi:A^{\rm gr}\xrightarrow{\sim}(S\otimes R,*_\pi),\qquad
 \phi\equiv1\pmod{\mathfrak m}
\end{equation}
が存在する。右辺も内部次数ごとの直和とする。
\end{lemma}
\begin{proof}
単位を入力すると零になるnormalized Hochschild cochainを使う。
多項式環の対角のgraded Koszul分解とnormalized bar分解の比較は
内部次数を保つように選べる。従ってpolydifferential cochainから
\emph{全ての}内部次数指定Hochschild cochainへの包含はHKR同型を誘導する。
既存の積がbidifferentialであるという仮定は不要である。
Kontsevichのグラフ写像はaffine線形変換に同変なので、内部weightを保つ。
weight零に制限したformalityと上の包含の合成は準同型であり、
次数別flatな結合積の変形をPoisson変形に対応させる。

この対応を各$R/\mathfrak m^{q+1}$で行う。
一段の核$\mathfrak m^q/\mathfrak m^{q+1}$では、持上げ障害の比較はHKRの
次数二での単射性、積の差の補正は次数一での全射性による。
一入力gaugeを同時に持ち上げるので、すでに選んだ低次jetを保てる。
これは\cite[Theorem~4.6]{PKontsevich}の冪零変形問題の比較の逐次形である。
全Hochschild cochain側のdegree零gaugeは一入力の写像であり、
厳密な代数同型を与える。normalized性から単位を保つ。
極限は各有限自由$A_j$内で存在し、逆写像も同様に存在する。
weight零のbivectorの係数次数は二である。
\end{proof}

\begin{lemma}[中心逆像を有限行列で取り出す]\label{cft:lem:center}
formalityを固定すると、$W\in S_3\otimes R$が一意に存在して
\begin{equation}\label{cft:eq:center}
 C_\pi(W)=\phi(\Omega),\qquad [\pi,W]=0,\qquad W\bmod\mathfrak m=F,
 \quad
 C_\pi(f):=\sum_{r\ge0}\frac1{r!}\mathcal U_{r+1}(f,\pi,\ldots,\pi)
\end{equation}
となる。
\end{lemma}
\begin{proof}
関数一個とbivector $r$個のグラフの出力arityは零であり、内部weightを保つ。
よって$C_\pi$は有限自由加群$S_3\otimes R$の$R$-線形自己写像で、
$\mathcal U_1(f)=f$により$C_\pi=1+O(\mathfrak m)$である。
その逆は各固定次数内で収束する幾何級数であり、$W$を一意に定める。

接写像のchain恒等式は
\begin{equation}\label{cft:eq:chain-center}
 [*_\pi,C_\pi(W)]=\mathcal U^\pi([\pi,W])
\end{equation}
である。左辺は中心性から零。
右辺の入力は係数次数四のvector場であり、接写像の最低次は
vector場を導分として見るHKRの単射である。
$\mathfrak m^q/\mathfrak m^{q+1}$で順にこの単射性を使うと
$[\pi,W]\in\bigcap_q\mathfrak m^q=0$となる。
曲率付きの通常のHochschild cohomologyをJacobian環と同一視する操作は使わない。
\end{proof}

\subsection{三次式の次数制約とゲージ経路}
\begin{lemma}[唯一の定数三重vectorと二回目の作用]\label{cft:lem:primitive}
一意な$\Xi\in\mathfrak m\bigwedge^3\langle\partial_i\rangle\otimes R$が存在し、
\begin{equation}\label{cft:eq:primitive}
 \pi=[\Xi,W],\qquad [\Xi,\pi]=0
\end{equation}
となる。
\end{lemma}
\begin{proof}
滑らかな$F$の全偏微分は多項式環の正則列である。
勾配Koszul contractionは、係数次数とrankの和「係数次数$+2\,\mathrm{rank}$」を保つ。
各斉次部分は有限の有限自由$R$-複体である。
剰余の完全性を使い、$\mathfrak m^q/\mathfrak m^{q+1}$で原像を逐次持ち上げると、
完備化と分離性から正のrankで完全となる。

$[\pi,W]=0$なので、係数次数二・rank二のcycle $\pi$の原像は
係数次数零・rank三に取れる。二つの原像の差の原像は係数次数$-2$を要求するため零。
これで$\Xi$の存在・一意性と$\Xi\bmod\mathfrak m=0$を得る。
$D=[\Xi,\pi]$は係数次数一・rank四であり、graded Jacobiと$[\pi,\pi]=0$から
$[W,D]=-[\pi,\pi]=0$となる。
そのKoszul原像は係数次数$-1$を要求するので$D=0$である。
ここでJacobi条件は必須であり、Casimir条件だけから第二の等式は従わない。
\end{proof}

周期変数$t$の二重次数を$(2,-3)$とし、$t$を微分せず、係数として多重線形に扱う。
$\eta=\Xi/t$はcohomological degree零・内部weight零である。
\begin{equation}\label{cft:eq:path}
 \alpha_s=-tW+(1-s)\pi,\qquad
 \partial_s\alpha_s=[\eta,\alpha_s]=-\pi
\end{equation}
はMaurer--Cartan経路で、始点は非可換模型、終点は可換模型に対応する。
$W$は$\mathfrak m$の倍数でないので、通常の冪零MC定理をこの経路へ直接適用してはいけない。

\begin{lemma}[大きい曲率でも有限なformality評価]\label{cft:lem:arity}
$\alpha_s$の像はcurved associative構造
$b_s=(*_{(1-s)\pi},-tC_{(1-s)\pi}(W))$である。
$\gamma_s=\mathcal U^{\alpha_s}(\eta)$はarity一と三だけを持ち、
$\gamma_s\in\mathfrak m$で、
\begin{equation}\label{cft:eq:lax}
 \partial_sb_s=[\gamma_s,b_s]
\end{equation}
となる。$\gamma_s$の正のarity成分はnormalizedである。
\end{lemma}
\begin{proof}
グラフの辺数と配置空間の次元$2n+m-2$を等しくすると、
関数$k$個とbivector任意個の出力arityは$2-2k$、
さらにtrivector一個を入れた場合は$3-2k$となる。
従って関数入力は高々一個で、残る入力は$\mathfrak m$の倍数である。
mod $\mathfrak m^{q+1}$では必要な全Taylor和が有限となり、
$L_\infty$恒等式も有限和の恒等式として適用できる。
chain恒等式に$\partial_s\alpha_s=-[\alpha_s,\eta]$を代入すると
$\partial_sb_s=-[b_s,\mathcal U^{\alpha_s}(\eta)]$、すなわち式\eqref{cft:eq:lax}である。

正のarityで単位をboundary入力へ置く。
そこへ辺が入れば定数の微分が零であり、入らなければ積分形式はその頂点を忘れた
一つ低次元の配置空間からの引戻しなので最高次積分が零となる。
internal頂点が一個の場合はHKR写像であり、同じ消滅は直接分かる。
この議論は\cite[Sections~6.2--6.4]{PKontsevich}のグラフ表示に適用する。
$t$への多重線形性と内部weightは全ての項で保たれる。
\end{proof}

\subsection{bar写像と逆写像の一様なarity評価}
\begin{lemma}[順序付きflowと逆写像]\label{cft:lem:flow}
$\partial_s\mathcal F_s=\gamma_s\mathcal F_s$、$\mathcal F_0=1$を解くと、
各$R/\mathfrak m^{q+1}$上で実際のbar coalgebra同型を得る。
$\mathcal F_1b_0=b_1\mathcal F_1$であり、写像・逆写像は式\eqref{cft:eq:bound}を満たす。
\end{lemma}
\begin{proof}
mod $\mathfrak m^{q+1}$で
\begin{equation}\label{cft:eq:ordered}
 \mathcal F_s=1+\sum_{j=1}^{q}
 \int_{0\le u_j\le\cdots\le u_1\le s}
 \gamma_{u_1}\cdots\gamma_{u_j}\,du_j\cdots du_1
\end{equation}
とする。各係数は$s$の多項式で、積分は標数零での代数的積分である。
$\gamma$がcoderivationなので、$\Delta\mathcal F_s$と
$(\mathcal F_s\otimes\mathcal F_s)\Delta$は同じ微分方程式・初期値を持つ。
よってcoalgebra写像である。
$\mathcal F_s-1\in\mathfrak m$から逆は
$\sum_{j=0}^{q}(1-\mathcal F_s)^j$で与えられる。
式\eqref{cft:eq:lax}により$\mathcal F_s^{-1}b_s\mathcal F_s$の微分は零である。

零項成分はなく、三入力coderivationを一回使うごとに出力長は二つ減る。
一入力成分は長さを変えない。従ってarity $2r+1$には少なくとも$r$回の
coderivationが必要で、係数は$\mathfrak m^r$に入る。
偶数arityは現れず、逆flowにも同じ評価が成立する。
さらに$t$を含まない入力に制限すれば
\begin{equation}\label{cft:eq:t-weight}
 F_{2r+1}=t^{-r}f_{2r+1},\qquad
 H_{2r+1}=t^{-r}g_{2r+1},
 \qquad \operatorname{wt}(f_{2r+1})=
 \operatorname{wt}(g_{2r+1})=-3r
\end{equation}
である。ここで$\operatorname{wt}$は、出力の多項式次数から入力の次数の和を引いた内部weightである。
実際、$\gamma_s$の三入力成分は関数入力を持たず、$\eta=\Xi/t$から$t^{-1}$を一つ含む。
一入力成分は関数入力$-tW$を一つ持つので$t$の冪は零である。
arityを$2r+1$から一へ減らすflowの項には三入力成分が正確に$r$回現れるため、
$t$の冪は$-r$である。内部weight零から$f_{2r+1},g_{2r+1}$のweightも従う。
出力tensor長$\ell$に寄与する成分の積の位数を足すと、mod $\mathfrak m^{q+1}$では
入力長は高々$\ell+2q$である。従ってtensor長完備barへの延長も各剰余で連続である。
単位はnormalized性からflow中固定され、高arity成分の単位入力は零となる。
同じ形式的係数を切断しているので、商$R/\mathfrak m^{q+1}\to R/\mathfrak m^q$と両立する。
\end{proof}

\subsection{冪零でないMF行列への評価と周期比較}
\begin{proof}[定理\ref{cft:thm:complete}の証明]
上のbar写像を有限個のshiftを持つ自由加群の行列成分に拡張する。
合成可能な添字列に沿う和とsuspensionのKoszul符号を用いるので、
元のbar恒等式は行列上でも成立する。
MF行列$Q$は$\mathfrak m$の倍数でも上三角でもなくてよい。
mod $\mathfrak m^{q+1}$では、像は有限式
\begin{equation}\label{cft:eq:object}
 Q'=\sum_{r=0}^{q}F_{2r+1}(Q,\ldots,Q)
\end{equation}
で与えられる。$k$個の射に対するTaylor成分は、suspended表示で
\begin{equation}\label{cft:eq:insertions}
 \sum_{r_0,\ldots,r_k\ge0}
 F_{k+r_0+\cdots+r_k}
 (Q_k^{\otimes r_k},a_k,Q_{k-1}^{\otimes r_{k-1}},\ldots,a_1,Q_0^{\otimes r_0})
\end{equation}
である。$k+\sum r_i\le2q+1$の項だけが寄与するため有限である。

suspensionを$\varsigma$と書くと、tensor長完備bar上で$\mathcal F(e^{\varsigma Q})=e^{\varsigma Q'}$である。
$\mathcal F b_0=b_1\mathcal F$から対象のMF方程式と全functor恒等式が従う。
同じ構成を逆bar写像に施すと、挿入公式はbar写像の合成と可換するので、
二つのtwisted functorの合成は恒等となる。
したがって各Artin切断で実際の互いに逆な関手がある。

有限rankの通常のgraded MFの形が保存されることも確認する。
基礎自由加群の偶・奇基底のcohomological次数をそれぞれ零・一に取ると、
その行列は、内部shiftを基底に付したまま
\begin{equation}\label{cft:eq:two-block-matrix}
 Q=\begin{pmatrix}0&tq_1\\q_0&0\end{pmatrix},
 \qquad Q^2=tC_\pi(W)\,\mathrm{id}
\end{equation}
と書ける。ここでは$\phi$で星積模型へ移した行列を用い、平方は始点の星積で計算する。
$2r+1$個の$Q$を挿入した合成可能な添字列では、
左下成分に$t$が$r$個、右上成分に$r+1$個現れる。
式\eqref{cft:eq:t-weight}の$t^{-r}$を掛けると、像の$t$の冪はそれぞれ零・一となる。
従って$Q'$は同じ有限個の内部shiftと同じ二つの階数を持つ行列であり、
bar恒等式により$(Q')^2=tW\,\mathrm{id}$を満たす。
逆写像についても同じ計算が成り立つ。

固定した有限自由加群、固定したcohomological次数と内部次数では、
各行列成分の$t$の冪と$x$の次数が固定される。
従って出力は有限個の$S_j\otimes R$の直和に属する。
式\eqref{cft:eq:bound}から、式\eqref{cft:eq:object}と\eqref{cft:eq:insertions}は
この有限自由加群内で$\mathfrak m$進収束する。
すべての恒等式は各Artin商で成立しているため、分離性から$R$上でも成立する。
これは同値の列の形式的な逆極限を宣言するのではなく、同じ基礎行列上の
関手と逆関手の明示的な係数wise極限である。
この極限には有限自由加群の$\mathfrak m$進完備性・分離性と、
各Artin商での有限評価だけを使っており、$R$の正則性や有限大域次元は必要ない。

内部shiftは基礎加群のラベルを変えるだけであり、行列上の公式は不変である。
各$F_r,H_r$は$R[t,t^{-1}]$-多重線形なので、$t$が与える
$(3)\simeq[2]$の自然比較とも可換する。
従って第\ref{gdr:sec:main}節で固定した周期比較がそのまま保存される。
正規化された単位も行列上で保たれる。
補題\ref{cft:lem:model}の厳密同型$\phi$と合成し、有限cone・冪等完備化へ延長して
式\eqref{cft:eq:equivalence}を得る。
\end{proof}
